\section{Simulation Results}
The simulation results demonstrating the system performance and functional correctness are shown in Fig.~\ref{fig:sim_metrics} and Fig.~\ref{fig:wave_fm}. The design implements a \textbf{fully pipelined FM Stereo Radio} receiver, achieving real-time streaming processing for a total of 32,768 output audio samples.

\subsection{Performance and Timing Metrics}

The performance metrics captured by the UVM monitor are summarized in Table \ref{tab:fm_perf}. The design achieves high efficiency with a deterministic latency between input streaming and high-fidelity audio output.

\begin{table}[htbp]
    \centering
    \caption{FM Radio Processor Performance Summary}
    \label{tab:fm_perf}
    \begin{tabular}{|l|r|}
        \hline
        \textbf{Metric} & \textbf{Value} \\
        \hline
        Total Audio Samples Captured & 32,768 \\
        \hline
        Clock Frequency & 100 MHz \\
        \hline
        First Valid Output Cycle & 71 \\
        \hline
        Total Simulation Cycles & 267,149 \\
        \hline
        Effective Throughput & 1 sample / cycle \\
        \hline
    \end{tabular}
\end{table}

The total simulation time of \textbf{267,149 cycles} for 32,768 output samples (at a decimation of 8) demonstrates the efficiency of the streaming FSM and the nested pipelining:
\begin{itemize}
    \item \textbf{High-throughput Input Buffer (\texttt{fifo.sv}):} A 16-deep FIFO is used to buffer USRP I/Q samples. The FSM manages data prefetching, matching the bursty nature of USRP data feeds while maintaining backpressure via the \texttt{in\_full} signal.
    \item \textbf{Pipelined DSP Pipeline:} Each stage (FIR, Demodulate, De-emphasis) operates over a multi-stage pipeline. The first valid audio sample appears at cycle 71, representing the intrinsic latency of the 32-stage divider and filtered register banks.
    \item \textbf{Deterministic Throughput:} Once the pipeline is full, the system produces one output sample every 8 input cycles (consistent with the \texttt{AUDIO\_DECIM} factor), utilizing the 100 MHz clock effectively.
\end{itemize}

\begin{figure}[htbp]
    \centering
    \includegraphics[width=1.0\linewidth]{Fig/waveform.png}
    \caption{Waveform of the FM Radio streaming data flow, showing stable FIFO operation and aligned Left/Right audio outputs.}
    \label{fig:wave_fm}
\end{figure}

Fig.~\ref{fig:wave_fm} provides a detailed waveform view of the processor's active computation:

\begin{itemize}
    \item \textbf{Streaming Data Input:} Under the \texttt{FM Radio Interfaces} group, data is streamed efficiently, with \texttt{valid\_in} and \texttt{in\_full} managing the flow control.
    \item \textbf{FIFO Management:} The \texttt{System Input FIFO} group displays internal count and read/write control signals, ensuring no data loss occurs despite the decimation stages downstream.
    \item \textbf{Audio Output:} The final outputs \texttt{left\_out} and \texttt{right\_out} are synchronized and bit-accurate to the reference software.
\end{itemize}

%--------------------------------------------------------------
\subsection{Bit-true comparison with software reference}
\begin{figure}[htbp]
    \centering
    \includegraphics[width=0.8\linewidth]{Fig/uvm_performance.png}
    \caption{UVM Performance Metrics and Scoreboard Results: 0 errors and 100\% Coverage.}
    \label{fig:sim_metrics}
\end{figure}

The verification was conducted using the \textbf{UVM Verification Environment (\texttt{my\_uvm\_tb.sv})}. Functionally validating all constraints via the \texttt{my\_uvm\_scoreboard.sv}, as evaluated in Fig.~\ref{fig:sim_metrics} and the UVM summary in Fig.~\ref{fig:uvm_summary}:
\begin{itemize}
    \item \textbf{Match Performance:} Both LEFT and RIGHT channels returned exactly \textbf{32,768 Checked} samples with \textbf{0 Errors}.
    \item \textbf{Result Cleanliness:} The UVM report summary confirmed exactly \textbf{0 UVM\_ERROR} and \textbf{0 UVM\_WARNING} reports across the entire dataset.
\end{itemize}
Bit-accuracy matches the C++ reference identically thanks to the 10-bit fixed-point quantization and the use of the \texttt{div1024\_f} function for consistent scaling.

\begin{figure}[htbp]
    \centering
    \includegraphics[width=0.4\linewidth]{Fig/uvm_report.png}
    \caption{UVM Report Summary showing a clean verification run with no errors.}
    \label{fig:uvm_summary}
\end{figure}


%--------------------------------------------------------------
\subsection{Functional coverage}

A UVM functional coverage model was explicitly implemented to check that the audio pipeline covers the full dynamic range of the synthesized signals.

\begin{table}[htbp]
    \centering
    \caption{Functional Coverage Results}
    \label{tab:coverage}
    \begin{tabular}{|c|p{6cm}|c|}
        \hline
        \textbf{Covergroup} & \textbf{Description} & \textbf{Coverage} \\
        \hline
        \texttt{cg\_audio\_output} & Validates output audio range (Positive, Negative, Zero) & 100.0\% \\
        \hline
    \end{tabular}
\end{table}

As shown in the simulation metrics (Fig.~\ref{fig:sim_metrics}), we attain \textbf{100.0\% Audio Output Coverage}. This proves that our test vectors successfully exercised the hardware across all boundary conditions, including positive and negative peaks in the audio waveform, ensuring the bit-true logic holds firm under all valid signal stimuli.

